\section{Experiments}
\label{sec:experiments}

\subsection{Experimental Setup}

\subsubsection{Hardware and Models}

Local retrieval benchmarks run on an Apple M-series machine (MPS/CPU, 16\,GB
unified memory); server experiments for ablation and comparison use an
NVIDIA RTX~3090 (24\,GB VRAM) with Ollama for LLM inference.
The LLM is Qwen3.5~9B~\citep{qwen2025} accessed via Ollama at
\texttt{localhost:11434}; the retrieval encoder is Snowflake
arctic-embed-m~\citep{merrick2024arctic} (109.5\,M parameters,
768-dimensional embeddings).
Neither model is fine-tuned.

\subsubsection{Evaluation Metrics}
\label{sec:metrics}

\paragraph{Execution Accuracy (EX).}
Let $\mathcal{Q} = \{q_1, \ldots, q_n\}$ denote the set of evaluation
questions, $G_i$ the gold SQL for question $q_i$, and $P_i$ the
predicted SQL.
Both are executed against the target SQLite database $d_i$;
$\text{exec}(S, d)$ denotes the result set (a multiset of rows)
returned by executing $S$ on $d$.
Execution accuracy is:
\begin{equation}
  \mathrm{EX} = \frac{1}{n}\sum_{i=1}^{n}
  \mathbf{1}\bigl[\text{exec}(P_i, d_i) = \text{exec}(G_i, d_i)\bigr].
  \label{eq:ex}
\end{equation}
Following \citet{gao2023dailsql}, result-set comparison is
column-order-insensitive and case-insensitive: row values are
normalised to lowercase strings and each row is represented as a
sorted tuple, so \texttt{SELECT a, b} and \texttt{SELECT b, a}
returning the same data are treated as equivalent.
If either the gold or predicted SQL raises a runtime error, the
question counts as incorrect (not excluded from the denominator);
queries where the gold SQL itself fails are excluded, accounting for
the four-question discrepancy between the 1{,}034-question dev.json
file and the 1{,}030 questions in the final evaluation
(four gold queries involve \texttt{CURDATE()} and are excluded as
unevaluable on offline SQLite).

\paragraph{Exact Match (EM).}
EM is a structural metric that ignores literal values:
\begin{equation}
  \mathrm{EM} = \frac{1}{n}\sum_{i=1}^{n}
  \mathbf{1}\bigl[\text{norm}(P_i) = \text{norm}(G_i)\bigr],
\end{equation}
where $\text{norm}(S)$ strips all string and numeric literals
(replacing them with \texttt{VALUE}) and normalises whitespace.
EM is a stricter metric than EX: a query that computes the
correct result via an algebraically equivalent but structurally
different expression (e.g.\ a self-join instead of a subquery)
counts as EX-correct but EM-incorrect.
EM serves as an upper bound on structural generalisation quality.

\paragraph{Recall at $K$ (R@$K$) for Schema Retrieval.}
Let $\mathcal{G}_q$ denote the set of gold tables for question $q$
(one or more tables depending on query complexity), and let
$\mathcal{R}_q^K$ denote the set of top-$K$ tables returned by the
retriever.
Retrieval success is defined under the lenient formulation:
\begin{equation}
  \mathrm{R@}K(q) = \mathbf{1}\bigl[
    \mathcal{G}_q \cap \mathcal{R}_q^K \neq \varnothing
  \bigr],
  \label{eq:rk}
\end{equation}
which counts a retrieval as a hit if at least one gold table appears
in the top-$K$ results.
The strict formulation would require $\mathcal{G}_q \subseteq \mathcal{R}_q^K$;
however, with a median of 2.1 gold tables per question on Spider~1.0
and $K$\,=\,5, the strict R@5 is 0.957 (FP32) versus 0.995 under the
lenient definition.\footnote{Lenient R@$K$ is reported throughout to remain
consistent with the DAIL-SQL evaluation protocol~\citep{gao2023dailsql}.}
The mean reciprocal rank is:
\begin{equation}
  \mathrm{MRR} = \frac{1}{n}\sum_{q=1}^{n}\frac{1}{\text{rank}_q},
  \label{eq:mrr}
\end{equation}
where $\text{rank}_q$ is the position of the first gold table in the
ranked retrieval list (or $\infty$ if no gold table appears, in which
case the term contributes~0).

\paragraph{Statistical Tests.}
All pairwise comparisons of two pipeline configurations use McNemar's
paired sign test with continuity correction.
For configurations $A$ and $B$, let $n_{01}$ be the number of questions
correct under $B$ but not $A$, and $n_{10}$ the reverse.
The test statistic is:
\begin{equation}
  \chi^2 = \frac{(|n_{01} - n_{10}| - 1)^2}{n_{01} + n_{10}},
\end{equation}
which follows a $\chi^2$ distribution with one degree of freedom under
the null hypothesis of equal accuracy.
Bootstrap 95\% confidence intervals for EX are computed using 10{,}000
resamples with replacement from the $n$~binary outcome values.

\subsection{Spider~1.0 Ablation Study}
\label{sec:exp_ablation}

Table~\ref{tab:ablation} reports execution accuracy for six incremental
pipeline configurations, evaluated on the 1{,}030 evaluable Spider~1.0
development questions.

\begin{table*}[!tbp]
  \centering
  \caption[Spider~1.0 pipeline ablation]{Execution accuracy (EX) and
    exact match (EM) for incremental pipeline configurations on
    1{,}034 Spider~1.0 development questions (fixed, column-order-insensitive
    evaluator).
    Each row adds one component over the previous best.
    All improvements over baseline are statistically significant
    ($p < 0.001$, McNemar's test).
    Errors denotes the number of SQL queries that raised a runtime
    execution error; $\Delta$EX is the gain over the baseline.}
  \label{tab:ablation}
  \begin{tabular}{lcccc}
    \toprule
    Configuration & EX & EM & Errors & $\Delta$EX \\
    \midrule
    Baseline ($k$=3, SC=1)     & 0.786 & 0.240 & 92 & ---       \\
    +sql2skeleton              & 0.805 & 0.284 & 46 & +1.91\,pp \\
    +$k$=5                     & 0.813 & 0.302 & 50 & +2.74\,pp \\
    +SC=3                      & 0.810 & 0.295 & 39 & +2.45\,pp \\
    +2-pass correction         & 0.823 & 0.345 & 47 & +3.71\,pp \\
    SC=3+$k$=5+2-pass (best)  & 0.824         & 0.344 & 38 & +3.87\,pp \\
    \bottomrule
  \end{tabular}
\end{table*}


Figure~\ref{fig:ablation} shows the incremental EX gain from each component.
The best configuration (SC\,=\,3, $k$\,=\,5, two-pass correction,
sql2skeleton) achieves EX\,=\,0.824 with a 95\% bootstrap confidence
interval of [0.800, 0.848].
Several aspects of Table~\ref{tab:ablation} require interpretation.

\begin{figure*}[!tbp]
  \centering
  \includegraphics[width=\linewidth]{fig2_ablation.pdf}
  \caption[Incremental pipeline ablation gains]{
    Incremental execution accuracy (EX) as each component is added to
    the baseline ($k$\,=\,3, SC\,=\,1).
    Each bar extends the running best; the final configuration
    (SC\,=\,3\,+\,$k$\,=\,5\,+\,2-pass) reaches EX\,=\,0.824.
    sql2skeleton delivers the largest single-component gain on Hard
    queries (+17.1\,pp); two-pass correction is the dominant contributor
    to overall EX (+3.71\,pp alone).
    All improvements are statistically significant ($p < 0.001$, McNemar's
    test with bootstrap 95\% CI).}
  \label{fig:ablation}
\end{figure*}

\paragraph{Non-monotonic errors column.}
The ``Errors'' column counts queries for which the final selected SQL
raises a SQLite runtime exception (e.g.\ \texttt{no such column},
\texttt{syntax error}).
This count is not monotonically decreasing across rows for two reasons.
First, the \emph{executor selector} in SC\,=\,3 preferentially returns the
first error-free candidate; when all three candidates fail, the first is
returned with its error exposed in the count.
Consequently, SC\,=\,3 can in some cases \emph{surface} errors that a
single-pass run would have silently produced as an incorrect result
(counted in EX as wrong, but not as an error).
Second, the two-pass corrector applies one re-prompt for the failing
query; this second attempt may produce a syntactically valid but
semantically incorrect SQL---in which case the error count decreases but EX
may not improve, explaining the +47 errors row versus +SC=3's 39.
The EX column (Eq.~\ref{eq:ex}) is the primary metric; the error count
is an auxiliary diagnostic of runtime failure modes.

\paragraph{EM--EX gap.}
The exact match rate (EM) is consistently lower than EX across all
configurations (e.g.\ EM\,=\,0.344 vs.\ EX\,=\,0.824 for the best
configuration), indicating that the model often produces algebraically
equivalent SQL via a different structural path than the gold query.
This is expected for a generation model with no task-specific training:
Qwen3.5~9B generates valid, result-equivalent SQL but does not reproduce
the specific join order, alias convention, or subquery structure of the
gold annotation.
The EM value should therefore not be interpreted as an accuracy penalty
but as a measure of structural alignment with the gold expression.

\paragraph{Synergy between SC=3 and two-pass correction.}
SC\,=\,3 alone adds +2.45\,pp; two-pass correction alone adds +3.71\,pp;
their combination adds +3.87\,pp over the sql2skeleton+$k$\,=\,5 baseline.
The sub-additive combination (+3.87\,pp~$<$~+2.45\,pp + 3.71\,pp)
arises because both components partially address the same population of
failures: SC\,=\,3 samples diverse structures to avoid systematic
errors, while two-pass correction handles the residual runtime failures.
The combination's marginal gain over two-pass alone (+0.16\,pp) confirms
that two-pass correction is the dominant single component, with SC\,=\,3
providing complementary diversity that enables the correction step to
succeed on previously unsalvageable candidates.

\subsection{Difficulty-Stratified Analysis}
\label{sec:exp_difficulty}

Table~\ref{tab:difficulty} and Figure~\ref{fig:difficulty} break down
EX by the four Spider difficulty levels defined by~\citet{yu2018spider}.
Difficulty is determined by SQL structural complexity: Easy queries
involve a single \texttt{SELECT} with at most one condition; Medium
queries add \texttt{JOIN} or \texttt{GROUP BY}; Hard queries use
multiple joins and \texttt{HAVING}; Extra Hard queries contain nested
subqueries (\texttt{NOT IN (SELECT\ldots)}, \texttt{INTERSECT},
\texttt{EXCEPT}).
The Spider~1.0 development set comprises 248 Easy, 446 Medium, 174 Hard,
and 166 Extra Hard questions.

\begin{table*}[!tbp]
  \centering
  \caption[EX by Spider difficulty level per pipeline configuration]{
    Execution accuracy by SQL difficulty level for each pipeline configuration.
    Distribution: Easy 248, Medium 446, Hard 174, Extra Hard 166
    (total 1{,}034).
    \textbf{Bold} marks the highest value per difficulty column.}
  \label{tab:difficulty}
  \begin{tabular}{lcccc}
    \toprule
    Configuration & Easy & Medium & Hard & Extra Hard \\
    \midrule
    Baseline ($k$=3, SC=1)      & 0.873 & 0.853 & 0.624 & 0.626 \\
    +sql2skeleton               & 0.874 & 0.865 & 0.755 & 0.574 \\
    +$k$=5                      & 0.880 & 0.870 & 0.773 & 0.585 \\
    +SC=3                       & 0.886 & 0.867 & 0.773 & 0.576 \\
    Best (SC=3+$k$=5+2pass)     & \textbf{0.893} & \textbf{0.884} & \textbf{0.795} & \textbf{0.587} \\
    \midrule
    $\Delta$ Baseline→Best      & +2.0\,pp & +3.1\,pp & +17.1\,pp & --3.9\,pp \\
    \bottomrule
  \end{tabular}
\end{table*}

\begin{figure*}[!tbp]
  \centering
  \includegraphics[width=\linewidth]{fig16_difficulty_breakdown.pdf}
  \caption[Difficulty-stratified EX: distribution, configuration comparison,
           and component gains]{
    (a)~The Spider~1.0 development set skews towards Medium queries
    (43\%), with Hard and Extra Hard each comprising 16\%.
    (b)~EX by difficulty for each pipeline configuration;
    Easy and Medium are near ceiling while Hard is the primary
    beneficiary of the full pipeline.
    (c)~Heatmap of $\Delta$EX (percentage points) for each added
    component versus the baseline; red cells indicate degradation,
    green cells indicate improvement.}
  \label{fig:difficulty}
\end{figure*}

Several patterns in Table~\ref{tab:difficulty} require explanation.

\paragraph{sql2skeleton: large Hard gain, Extra Hard degradation.}
Adding sql2skeleton raises Hard EX by +13.1\,pp (0.624$\to$0.755)
but \emph{lowers} Extra Hard EX by $-$5.2\,pp (0.626$\to$0.574).
Hard queries benefit because their structural skeleton (multi-join +
\texttt{HAVING count(*) >= VALUE}) is distinctive enough to retrieve
few-shot examples with an identical structural template, directly
supplying the LLM with a worked JOIN+HAVING demonstration.
Extra Hard queries, however, often involve rare nested operators
(\texttt{INTERSECT}, \texttt{EXCEPT}, correlated subqueries) for which
the training pool contains few exact structural matches; skeleton-based
retrieval in this case returns less relevant examples than question-text
similarity, which can still locate domain-relevant content.
This tradeoff explains why the overall gain from sql2skeleton (+1.91\,pp)
is dominated by the Hard category (+13.1\,pp $\times$ 174/1034 normalised
contribution = +2.2\,pp) partially offset by the Extra Hard loss.

\paragraph{Extra Hard ceiling.}
Even the best configuration reaches only 0.587 on Extra Hard, a
26.1\,pp gap below the Easy ceiling of 0.893.
Extra Hard queries require correctly placing nested subqueries, correlated
references, and set operators (\texttt{INTERSECT}/\texttt{UNION}) that a
9B-parameter model with no SQL fine-tuning handles unreliably.
The two-pass corrector recovers some cases (runtime errors surface
immediately as feedback), but semantic errors in nested queries
(wrong correlated condition) are not detectable from execution errors alone
and remain uncorrected.

\paragraph{Easy--Medium near ceiling.}
The full pipeline achieves 0.893 on Easy and 0.884 on Medium, indicating
that straightforward single-table and single-join queries are reliably
handled by the LLM with minimal few-shot guidance.
The remaining gap from 1.0 is attributable to ambiguous natural-language
phrasings (e.g.\ ordinal vs.\ cardinal count) and occasional schema
mismatch between retrieved tables and the gold query.

\subsection{Comparison with Published Systems}

Table~\ref{tab:sota} situates the pipeline relative to published
Text-to-SQL systems on Spider~1.0.

\begin{table*}[!tbp]
  \centering
  \caption[Comparison with published Text-to-SQL systems]{Execution
    accuracy on Spider~1.0 development set for representative published
    systems.
    Systems are grouped by base model type.
    Our pipeline uses a 9B-parameter local model with no fine-tuning.
    DAIL-SQL and DIN-SQL use GPT-4 ($\sim$1.7\,T parameters) via
    cloud API; RESDSQL and PICARD use fine-tuned 3B-parameter T5 models.
    The comparison demonstrates that the proposed pipeline is competitive
    with fine-tuned 3B-parameter models while requiring no training and
    running entirely on consumer hardware.}
  \label{tab:sota}
  \begin{tabular}{lllc}
    \toprule
    System & Base model & Training & EX \\
    \midrule
    DAIL-SQL~\citep{gao2023dailsql}     & GPT-4 ($\sim$1.7\,T) & few-shot & 0.866 \\
    DIN-SQL~\citep{pourreza2023dinsql}  & GPT-4 ($\sim$1.7\,T) & few-shot & 0.826 \\
    \midrule
    RESDSQL~\citep{li2023resdsql}       & T5-3B               & fine-tuned & 0.791 \\
    PICARD~\citep{scholak2021picard}    & T5-3B               & fine-tuned & 0.754 \\
    \midrule
    Ours (FP32 retriever)$^\dagger$    & Qwen3.5 9B (local)  & none       & 0.824 \\
    \bottomrule
  \end{tabular}
  \smallskip
  {\footnotesize $^\dagger$\,Retrieval quality under INT8 is identical to
  FP32 ($\Delta$R@5\,=\,0; Section~\ref{sec:quantization}); a paired
  downstream comparison on identical hardware is left for future work.}
\end{table*}

The pipeline achieves EX\,=\,0.824, which is competitive with fine-tuned
3B-parameter models (RESDSQL: 0.791, PICARD: 0.754) and with DIN-SQL
(0.826), which uses a model approximately 190$\times$ larger.

The 4.2\,pp gap relative to DAIL-SQL (0.866) has two identifiable sources.
First, GPT-4 has an estimated 1.7\,T parameters versus 9B here---a 190$\times$
difference in model capacity that enables more reliable multi-hop reasoning
for extra-hard queries (nested subqueries, correlated aggregates).
Second, DAIL-SQL uses the full DAIL selection pool with cross-DB few-shot
examples from 147 training databases, whereas this work uses a 24-example
Spider-only pool, providing fewer structural templates.\
Both sources are capacity constraints of the edge deployment setting,
not deficiencies of the pipeline design itself.
The DIN-SQL comparison (0.826 vs.\ 0.824, $\Delta$\,=\,$-$0.002) is not
statistically significant at $\alpha\,=\,0.05$ and should be interpreted
as parity.

\subsection{Spider~2.0-Lite Results}

Table~\ref{tab:sp2} reports execution accuracy on the Spider~2.0-Lite
local SQLite subset (135 questions).

\begin{table}[!tbp]
  \centering
  \caption[Spider~2.0-Lite execution accuracy]{Execution accuracy (EX)
    on the Spider~2.0-Lite local SQLite subset (135 questions across
    30 databases) for three pipeline configurations.
    Zero-shot refers to direct prompting without few-shot examples.
    Few-shot uses $k$\,=\,3 examples retrieved from a 24-example gold
    SQL pool via arctic-embed-m.
    SC\,=\,3 + 2-pass denotes the full pipeline.
    The large gain from zero-shot to few-shot (+3.7\,pp) and the
    substantially lower absolute accuracy compared to Spider~1.0
    reflect the analytical complexity of Spider~2.0-Lite queries.}
  \label{tab:sp2}
  \begin{tabular}{lcc}
    \toprule
    Configuration & EX & Correct/Total \\
    \midrule
    Zero-shot            & 0.067 & 9/135  \\
    +few-shot $k$=3      & 0.104 & 14/135 \\
    +SC=3 + 2-pass       & \textbf{0.111}   & \textbf{15/135} \\
    \bottomrule
  \end{tabular}
\end{table}

Each pipeline component contributes incrementally on Spider~2.0-Lite:
few-shot retrieval raises EX from 0.067 to 0.104 (+55\% relative),
and adding SC\,=\,3 with two-pass correction yields a further gain to
0.111 (+7\% relative, 15 correct queries).

\paragraph{Analysis of the low absolute accuracy on Spider~2.0-Lite.}
The EX\,=\,0.111 result is low relative to Spider~1.0 (0.824) and
requires explanation rather than dismissal.
Error analysis (Appendix~\ref{app:error_taxonomy})
identifies three structural root causes.
First, \texttt{result\_mismatch} (63 queries) accounts for the largest
failure category: the predicted SQL executes without error but returns
an incorrect result set.
These failures arise from semantically complex constructs (window functions,
\texttt{RANK() OVER PARTITION BY}, recursive CTEs) that appear in
Spider~2.0-Lite but are essentially absent from Spider~1.0, leaving
the 9B model without structural templates from the few-shot pool.
Second, \texttt{no\_such\_column} (33 queries) reflects schema ambiguity:
Spider~2.0-Lite uses non-standard column names (e.g.\ \texttt{cdscode},
\texttt{avgscread}) without the descriptive conventions Spider~1.0 relies on,
reducing the signal available to the retriever and the LLM.
Third, the few-shot pool covers only 16 of 30 databases; 34\% of test
questions have no same-database example, eliminating the primary mechanism
through which few-shot retrieval reduces \texttt{no\_such\_column} errors.
Despite these limitations, the consistent incremental gain from each
component ($+55\%$ from few-shot, $+7\%$ from SC\,=\,3+2-pass) confirms
that the pipeline design transfers to harder benchmarks even when
the absolute ceiling is constrained by model capacity.

\subsection{Retrieval Quality Comparison: FP32 vs.\ INT8}

Figure~\ref{fig:rk_comparison} provides a direct visual comparison of
FP32 and INT8 R@K on both Spider~1.0 (1{,}034 queries, schema-linking
ground truth) and Spider~2.0-Lite (89 evaluable queries, DB-match
ground truth).

\begin{figure*}[!tbp]
  \centering
  \includegraphics[width=\linewidth]{fig5_rk_comparison.pdf}
  \caption[R@K comparison: FP32 vs.\ INT8 on Spider~1.0 and 2.0-Lite]{
    Recall at $K$ ($K \in \{1,3,5,10\}$) for FP32 and INT8 retrieval
    profiles on Spider~1.0 (left) and Spider~2.0-Lite (right).
    On Spider~1.0, INT8 matches FP32 exactly at all $K$ values
    ($\Delta$R@K\,=\,0.0000 for all $K$).
    On Spider~2.0-Lite, INT8 R@10 is identical to FP32; the
    $\Delta$R@5\,=\,$-$0.023 corresponds to two queries out of 89
    and is not statistically significant given the small pool size.}
  \label{fig:rk_comparison}
\end{figure*}

